% =====================================================================
\chapter{Integrales impropias}
\label{tema:4}
% =====================================================================

La integral de Riemann exige un intervalo cerrado y acotado y una función acotada. Muchas integrales importantes no cumplen alguna de las dos condiciones: la densidad normal está definida en toda la recta real, la función Gamma integra en $[0,\infty)$ y la densidad de una ley de potencia no está acotada cerca del origen si el exponente es negativo. Este tema extiende la integral mediante un paso al límite, estudia cuándo ese límite existe y calcula dos integrales fundamentales para la probabilidad: la función Gamma y la integral gaussiana. Se siguen la §3.7 de~\cite{openstax2}, la §7.5 de Strang~\cite{strang}, el capítulo~10 de Apostol~\cite{apostol} y las sesiones~91 a~93 del curso 18.01SC~\cite{mit1801}.

\section{Integrales en intervalos no acotados}

\begin{definicion}[Integral impropia de primera especie]\label{def:impropia1}
Sea $f\colon[a,\infty)\to\R$ integrable en $[a,b]$ para todo $b>a$. Si existe el límite finito
\[
\int_a^{\infty}f(x)\dif x=\lim_{b\to\infty}\int_a^{b}f(x)\dif x,
\]
se dice que la integral impropia \emph{converge} y que ese es su valor; en otro caso, \emph{diverge}. Se define igual $\int_{-\infty}^{b}f$. Para una función definida en todo $\R$, se dice que $\int_{-\infty}^{\infty}f$ converge si convergen $\int_{-\infty}^{c}f$ y $\int_{c}^{\infty}f$ para algún $c$, y su valor es la suma de ambas.
\end{definicion}

Por la aditividad, la elección de $c$ no afecta ni a la convergencia ni al valor. La definición exige que converjan las dos colas por separado: $\int_{-\infty}^{\infty}x\dif x$ diverge, aunque $\int_{-R}^{R}x\dif x=0$ para todo $R$.

\begin{proposicion}[Integrales $p$ en el infinito]\label{prop:p-infinito}
Sea $p\in\R$. La integral $\int_1^{\infty}\frac{\dif x}{x^{p}}$ converge si y solo si $p>1$, y en ese caso vale $\frac{1}{p-1}$.
\end{proposicion}

\begin{proof}
Si $p\neq1$, $\int_1^{b}x^{-p}\dif x=\frac{b^{1-p}-1}{1-p}$. Si $p>1$, $b^{1-p}\to0$ y el límite es $\frac{-1}{1-p}=\frac{1}{p-1}$. Si $p<1$, $b^{1-p}\to\infty$ y la integral diverge. Si $p=1$, $\int_1^{b}\frac{\dif x}{x}=\ln b\to\infty$.
\end{proof}

\begin{ejemplo}\label{ej:impropias1}
(a) Para $a>0$, $\int_0^{\infty}e^{-ax}\dif x=\lim_{b\to\infty}\frac{1-e^{-ab}}{a}=\frac1a$.
(b) $\int_{-\infty}^{\infty}\frac{\dif x}{1+x^{2}}=\lim_{b\to\infty}\arctan b-\lim_{a\to-\infty}\arctan a=\frac\pi2-\bigl(-\frac\pi2\bigr)=\pi$.
\end{ejemplo}

\section{Integrandos no acotados}

\begin{definicion}[Integral impropia de segunda especie]\label{def:impropia2}
Sea $f\colon(a,b]\to\R$ integrable en $[c,b]$ para todo $c\in(a,b)$ y no acotada cerca de $a$. Se define $\int_a^{b}f=\lim_{c\to a^{+}}\int_c^{b}f$ si el límite existe y es finito. Análogamente si la singularidad está en $b$. Si hay una singularidad en un punto interior, o varias, se divide el intervalo y se exige que converja cada trozo.
\end{definicion}

\begin{proposicion}[Integrales $p$ en el origen]\label{prop:p-cero}
La integral $\int_0^{1}\frac{\dif x}{x^{p}}$, con $p>0$, converge si y solo si $p<1$, y en ese caso vale $\frac{1}{1-p}$.
\end{proposicion}

\begin{proof}
Si $p\neq1$, $\int_c^{1}x^{-p}\dif x=\frac{1-c^{1-p}}{1-p}$. Si $p<1$, $c^{1-p}\to0$ cuando $c\to0^{+}$ y el límite es $\frac{1}{1-p}$; si $p>1$, $c^{1-p}\to\infty$. Si $p=1$, $\int_c^{1}\frac{\dif x}{x}=-\ln c\to\infty$.
\end{proof}

Las \cref{prop:p-infinito,prop:p-cero} son simétricas: $x^{-p}$ es integrable en el infinito si decrece más deprisa que $1/x$, y en el origen si crece más despacio que $1/x$. En particular, $\int_0^{\infty}x^{-p}\dif x$ diverge para todo $p$.

\begin{ejemplo}\label{ej:log-impropia}
$\int_0^{1}\ln x\dif x$: por el \cref{ej:partes}(b), $\int_c^{1}\ln x\dif x=\eval{x\ln x-x}{c}{1}=-1-c\ln c+c$. Como $c\ln c\to0$ cuando $c\to0^{+}$ (regla de L'Hôpital, módulo~0.5), la integral converge y vale $-1$.
\end{ejemplo}

\section{Criterios de convergencia}

Casi nunca se conoce una primitiva explícita, y conviene decidir la convergencia comparando con integrales conocidas. Los criterios se enuncian en $[a,\infty)$; para las de segunda especie son idénticos.

\begin{teorema}[Criterio de comparación]\label{teo:comparacion-int}
Sean $f,g$ integrables en $[a,b]$ para todo $b>a$, con $0\le f(x)\le g(x)$ para $x\ge a$. Si $\int_a^{\infty}g$ converge, $\int_a^{\infty}f$ converge y $\int_a^{\infty}f\le\int_a^{\infty}g$. Si $\int_a^{\infty}f$ diverge, $\int_a^{\infty}g$ diverge.
\end{teorema}

\begin{proof}
Sea $F(b)=\int_a^{b}f$. Como $f\ge0$, $F$ es creciente, y por la monotonía de la integral $F(b)\le\int_a^{b}g\le\int_a^{\infty}g$, pues también $b\mapsto\int_a^b g$ es creciente con límite $\int_a^\infty g$. Sea $S=\sup_{b>a}F(b)$, que existe porque $F$ está acotada. Dado $\varepsilon>0$, hay $b_0$ con $F(b_0)>S-\varepsilon$, y por la monotonía $S-\varepsilon<F(b)\le S$ para todo $b\ge b_0$. Luego $F(b)\to S$. La segunda afirmación es la contrarrecíproca de la primera.
\end{proof}

\begin{teorema}[Comparación en el límite]\label{teo:comparacion-limite}
Sean $f,g>0$ en $[a,\infty)$, integrables en cada $[a,b]$, con $\lim_{x\to\infty}\frac{f(x)}{g(x)}=L\in(0,\infty)$. Entonces $\int_a^{\infty}f$ y $\int_a^{\infty}g$ son ambas convergentes o ambas divergentes.
\end{teorema}

La demostración se pide en el \cref{eje:4-limite-comp}. En la práctica se compara con $x^{-p}$: por ejemplo, $\frac{1}{\sqrt{x^{3}+1}}$ se comporta como $x^{-3/2}$ en el infinito, y $\int_1^{\infty}\frac{\dif x}{\sqrt{x^{3}+1}}$ converge.

\begin{teorema}[Convergencia absoluta]\label{teo:absoluta-int}
Sea $f$ integrable en cada $[a,b]$. Si $\int_a^{\infty}\abs{f}$ converge, también converge $\int_a^{\infty}f$, y $\bigl|\int_a^{\infty}f\bigr|\le\int_a^{\infty}\abs{f}$.
\end{teorema}

\begin{proof}
Se cumple $0\le f+\abs{f}\le2\abs{f}$. Por el \cref{teo:comparacion-int}, $\int_a^{\infty}(f+\abs{f})$ converge. Entonces $\int_a^{b}f=\int_a^{b}(f+\abs{f})-\int_a^{b}\abs{f}$ es diferencia de dos funciones de $b$ con límite finito, y tiene límite. La desigualdad se obtiene pasando al límite en $\bigl|\int_a^{b}f\bigr|\le\int_a^{b}\abs{f}$.
\end{proof}

\begin{ejemplo}\label{ej:criterios}
(a) Para $x\ge1$, $x^{2}\ge x$, luego $0<e^{-x^{2}}\le e^{-x}$, y $\int_1^{\infty}e^{-x^{2}}\dif x$ converge por comparación con el \cref{ej:impropias1}(a).
(b) $\bigl|\frac{\sin x}{x^{2}}\bigr|\le\frac{1}{x^{2}}$, luego $\int_1^{\infty}\frac{\sin x}{x^{2}}\dif x$ converge absolutamente y, por el \cref{teo:absoluta-int}, converge.
(c) $\int_1^{\infty}\frac{\dif x}{x+\sqrt x}$ diverge: $\frac{1/(x+\sqrt x)}{1/x}=\frac{1}{1+x^{-1/2}}\to1$ y $\int_1^{\infty}\frac{\dif x}{x}$ diverge.
\end{ejemplo}

\section{La función Gamma}

Para todo $n\in\N_0$, el límite $\lim_{x\to\infty}x^{n}e^{-x}=0$ (módulo~0.5) expresa que la exponencial domina a cualquier potencia. Gracias a ello las integrales $\int_0^{\infty}x^{n}e^{-x}\dif x$ convergen.

\begin{definicion}[Función Gamma]\label{def:gamma}
Para $s>0$ se define $\Gamma(s)=\int_0^{\infty}x^{s-1}e^{-x}\dif x$.
\end{definicion}

La integral es impropia en el infinito y, si $s<1$, también en el origen. Converge para todo $s>0$: en $(0,1]$, $0<x^{s-1}e^{-x}\le x^{s-1}=x^{-(1-s)}$ con $1-s<1$ (\cref{prop:p-cero}); en $[1,\infty)$, como $x^{s+1}e^{-x}\to0$, existe $x_0$ con $x^{s-1}e^{-x}\le x^{-2}$ para $x\ge x_0$ (\cref{prop:p-infinito}). El \cref{eje:4-gamma-conv} prueba que diverge si $s\le0$.

\begin{proposicion}\label{prop:gamma-factorial}
Para todo $s>0$, $\Gamma(s+1)=s\,\Gamma(s)$. En particular, $\Gamma(1)=1$ y $\Gamma(n+1)=n!$ para todo $n\in\N_0$.
\end{proposicion}

\begin{proof}
Sean $0<c<b$. Por partes con $u=x^{s}$, $v'=e^{-x}$:
$\int_c^{b}x^{s}e^{-x}\dif x=\eval{-x^{s}e^{-x}}{c}{b}+s\int_c^{b}x^{s-1}e^{-x}\dif x$.
Cuando $b\to\infty$ y $c\to0^{+}$, el término de borde tiende a $0$, porque $b^{s}e^{-b}\to0$ y $c^{s}e^{-c}\to0$ al ser $s>0$; queda $\Gamma(s+1)=s\Gamma(s)$. Además $\Gamma(1)=\int_0^{\infty}e^{-x}\dif x=1$, y por inducción $\Gamma(n+1)=n\,\Gamma(n)=n\,(n-1)!=n!$.
\end{proof}

La función Gamma extiende el factorial a valores no enteros; por ejemplo, $\Gamma(\frac12)=\sqrt\pi$ (\cref{eje:4-gamma-medio}).

\section{La integral gaussiana}

La función $e^{-x^{2}}$ no tiene primitiva elemental, pero su integral sobre toda la recta puede calcularse exactamente. Es la constante que normaliza la distribución normal.

\begin{teorema}[Integral gaussiana]\label{teo:gauss}
$\displaystyle\int_{-\infty}^{\infty}e^{-x^{2}}\dif x=\sqrt\pi$.
\end{teorema}

\begin{proof}
\emph{Convergencia.} Para $\abs{x}\ge1$, $e^{-x^{2}}\le e^{-\abs{x}}$, y el \cref{teo:comparacion-int} con el \cref{ej:impropias1}(a) prueba que convergen las dos colas. Sea $I$ el valor; como el integrando es par, $I=2\int_0^\infty e^{-x^2}\dif x>0$.

\emph{Cálculo.} El cuadrado de $I$ es una integral doble sobre el plano:
\[
I^{2}=\int_{-\infty}^{\infty}e^{-x^{2}}\dif x\int_{-\infty}^{\infty}e^{-y^{2}}\dif y=\iint_{\R^{2}}e^{-(x^{2}+y^{2})}\dif x\dif y .
\]
El integrando depende solo de la distancia $r=\sqrt{x^{2}+y^{2}}$ al origen. Se descompone el plano en coronas circulares de radio $r$ y grosor $\dif r$: el área de la corona entre $r$ y $r+\dif r$ es $\pi(r+\dif r)^{2}-\pi r^{2}\approx2\pi r\dif r$, y en ella el integrando vale aproximadamente $e^{-r^{2}}$. Sumando las coronas,
\[
I^{2}=\int_0^{\infty}e^{-r^{2}}\,2\pi r\dif r=\pi\lim_{R\to\infty}\eval{-e^{-r^{2}}}{0}{R}=\pi .
\]
Como $I>0$, $I=\sqrt\pi$.
\end{proof}

\begin{observacion}[Sobre el paso a coordenadas polares]
La demostración usa dos hechos del cálculo en varias variables: que el producto de dos integrales es una integral doble (teorema de Fubini) y la fórmula del cambio a coordenadas polares, $\dif x\dif y=r\dif r\dif\theta$, de la que el argumento de las coronas es la versión intuitiva. Ambos se justifican en el módulo~0.7; véase también la §14.2 «Changing to Better Coordinates» de Strang~\cite{strang}. El \cref{eje:4-wallis-gauss} obtiene el mismo resultado sin salir de una variable, con las integrales de Wallis del \cref{eje:3-wallis}.
\end{observacion}

\begin{corolario}\label{cor:gauss}
Para $\mu\in\R$ y $\sigma>0$,
\[
\int_{-\infty}^{\infty}e^{-x^{2}/2}\dif x=\sqrt{2\pi},\qquad\int_{-\infty}^{\infty}e^{-(x-\mu)^{2}/(2\sigma^{2})}\dif x=\sigma\sqrt{2\pi}.
\]
En consecuencia, $\varphi(x)=\frac{1}{\sigma\sqrt{2\pi}}e^{-(x-\mu)^{2}/(2\sigma^{2})}$ es una densidad de probabilidad: es positiva e integra~1.
\end{corolario}

\begin{proof}
En cada integral sobre $[a,b]$ se hace el cambio lineal $x=\mu+\sigma\sqrt2\,t$, $\dif x=\sigma\sqrt2\dif t$ (\cref{teo:sustitucion}), y se pasa al límite: se obtiene $\sigma\sqrt2\int_{-\infty}^{\infty}e^{-t^{2}}\dif t=\sigma\sqrt{2\pi}$. La primera igualdad es el caso $\mu=0$, $\sigma=1$.
\end{proof}

\begin{observacion}[Errores frecuentes]\label{obs:4-errores}
(a) No detectar que una integral es impropia porque el integrando se anula o explota en un punto interior, como $\int_{-1}^{1}\frac{\dif x}{x^{2}}$. (b) Calcular $\int_{-\infty}^{\infty}f$ como $\lim_{R\to\infty}\int_{-R}^{R}f$; ese límite (el \emph{valor principal}) puede existir aunque la integral diverja. (c) Creer que $f(x)\to0$ basta para que $\int_a^{\infty}f$ converja: $\frac1x\to0$ y la integral diverge. (d) Aplicar el criterio de comparación a funciones que cambian de signo; en ese caso se compara $\abs{f}$ y se usa el \cref{teo:absoluta-int}.
\end{observacion}

\begin{observacion}[Conexión con el curso]
En el módulo~5, toda densidad cumple $\int_{-\infty}^{\infty}f_X=1$, una integral impropia, y el \cref{cor:gauss} es la normalización de la distribución normal; la varianza de la normal estándar es $\int x^{2}\varphi(x)\dif x=1$ (\cref{eje:4-normal}). En el módulo~11, los momentos de una activación aplicada a una preactivación gaussiana, como $\E[\mathrm{ReLU}(z)^{2}]$, son integrales de este tipo. En el módulo~16, una variable con densidad proporcional a $x^{-\alpha}$ en $[1,\infty)$, una ley de potencia, solo tiene media finita si $\alpha>2$, consecuencia directa de la \cref{prop:p-infinito} (\cref{eje:4-potencias}).
\end{observacion}

\begin{lecturas}
\item OpenStax~\cite{openstax2}, §3.7 «Improper Integrals».
\item Strang~\cite{strang}, §7.5 «Improper Integrals» y §8.4 «Probability and Calculus».
\item Apostol~\cite{apostol}, cap.~10 «Sequences, Infinite Series, and Improper Integrals» (integrales impropias).
\item MIT 18.01SC~\cite{mit1801}, unidad~5, parte~A, sesiones~91 «Improper Integrals», 92 «Integral Comparison» y~93 «Indefinite Integrals and Singularities»; unidad~3, sesión~65 «Bell Curve, Conclusion».
\end{lecturas}

\section*{Ejercicios del tema 4}
\addcontentsline{toc}{section}{Ejercicios del tema 4}

\begin{ejercicio}{1}{eje:4-clasificar}
Indique cuáles de las integrales son impropias y de qué especie (y en qué punto): (a) $\int_0^{1}\frac{\dif x}{\sqrt x}$; (b) $\int_0^{2}\frac{\dif x}{x-1}$; (c) $\int_1^{\infty}e^{-x}\dif x$; (d) $\int_0^{\pi/2}\tan x\dif x$; (e) $\int_0^{1}\frac{\sin x}{x}\dif x$; (f) $\int_{-\infty}^{\infty}\frac{\dif x}{1+x^{2}}$.
\end{ejercicio}

\begin{ejercicio}{2}{eje:4-calcular}
Calcule o demuestre que divergen: (a) $\int_1^{\infty}\frac{\dif x}{x^{2}}$; (b) $\int_1^{\infty}x^{-10}\dif x$; (c) $\int_0^{1}\frac{\dif x}{x}$; (d) $\int_0^{1}\frac{\dif x}{\sqrt x}$; (e) $\int_0^{e}\ln x\dif x$; (f) $\int_0^{\infty}x\,e^{-x^{2}}\dif x$.
\fuente{adaptado de Strang~\cite{strang}, §7.5, ejercicios 1, 2 y~9, y de MIT 18.01SC~\cite{mit1801}, \emph{18.01 Exercises}, ejercicio 3C-4.}
\end{ejercicio}

\begin{ejercicio}{2}{eje:4-logaritmos}
Estudie la convergencia, calculando cuando sea posible: (a) $\int_2^{\infty}\frac{\dif x}{x\ln x}$; (b) $\int_2^{\infty}\frac{\dif x}{x(\ln x)^{2}}$; (c) $\int_{100}^{\infty}\frac{\dif x}{x\,\ln x\,\ln(\ln x)}$. ¿Qué enseña (c) sobre la frontera entre convergencia y divergencia?
\fuente{adaptado de Strang~\cite{strang}, §7.5, ejercicio~11.}
\end{ejercicio}

\begin{ejercicio}{2}{eje:4-comparar}
Decida si convergen, comparando con una integral conocida: (a) $\int_1^{\infty}\frac{\dif x}{x^{3}+1}$; (b) $\int_1^{\infty}\frac{2+\sin x}{x}\dif x$; (c) $\int_0^{1}\frac{\dif x}{\sqrt{x+x^{3}}}$; (d) $\int_1^{\infty}\frac{\dif x}{\sqrt{x^{2}+x}}$; (e) $\int_0^{1}\frac{e^{x}}{x}\dif x$; (f) $\int_0^{\infty}\frac{x^{2}}{e^{x}}\dif x$.
\fuente{adaptado de Strang~\cite{strang}, §7.5, ejercicios 17--26.}
\end{ejercicio}

\begin{ejercicio}{2}{eje:4-seno}
Demuestre que $\int_1^{N}\frac{\sin x}{x^{2}}\dif x<1$ para todo $N>1$ y que $\int_1^{\infty}\frac{\sin x}{x^{2}}\dif x$ converge.
\fuente{adaptado de MIT 18.01SC~\cite{mit1801}, \emph{18.01 Exercises}, ejercicio 3E-7.}
\end{ejercicio}

\begin{ejercicio}{2}{eje:4-gamma-valores}
Con la \cref{prop:gamma-factorial}, calcule $\Gamma(2)$, $\Gamma(5)$ y $\frac{\Gamma(7)}{\Gamma(5)}$. Con el cambio $u=2x$, calcule $\int_0^{\infty}x^{3}e^{-2x}\dif x$ y, en general, $\int_0^{\infty}x^{n}e^{-\lambda x}\dif x$ para $\lambda>0$.
\end{ejercicio}

\begin{ejercicio}{2}{eje:4-gamma-medio}
(a) Con el cambio $x=u^{2}$, demuestre que $\Gamma(\frac12)=2\int_0^{\infty}e^{-u^{2}}\dif u=\sqrt\pi$. (b) Calcule $\Gamma(\frac32)$ y $\Gamma(\frac52)$; en la notación $s!=\Gamma(s+1)$, compruebe que $(\frac12)!=\frac{\sqrt\pi}{2}$.
\fuente{adaptado de Strang~\cite{strang}, §7.5, ejercicios 27 y~28.}
\end{ejercicio}

\begin{ejercicio}{2}{eje:4-erf}
Sea $E(x)=\frac{1}{\sqrt{2\pi}}\int_0^{x}e^{-u^{2}/2}\dif u$ y $F(x)=\int_0^{x}e^{-t^{2}}\dif t$. (a) Exprese $E$ en función de $F$ con un cambio de variable y, sabiendo que $F(x)\to\frac{\sqrt\pi}{2}$ cuando $x\to\infty$, calcule $\lim_{x\to\infty}E(x)$ y $\lim_{x\to-\infty}E(x)$. (b) Exprese $\frac{1}{\sqrt{2\pi}}\int_a^{b}e^{-u^{2}/2}\dif u$ con $E$.
\fuente{adaptado de MIT 18.01SC~\cite{mit1801}, \emph{18.01 Exercises}, ejercicio 3E-2.}
\end{ejercicio}

\begin{ejercicio}{2}{eje:4-normal}
(a) Integrando por partes con $u=x$ y $v'=x\,e^{-x^{2}/2}$, demuestre que $\int_{-\infty}^{\infty}x^{2}e^{-x^{2}/2}\dif x=\sqrt{2\pi}$, y deduzca que la normal estándar tiene varianza~1. (b) Calcule $\int_0^{\infty}x^{2}e^{-x^{2}}\dif x$.
\end{ejercicio}

\begin{ejercicio}{3}{eje:4-gamma-conv}
Demuestre que $\int_0^{\infty}x^{s-1}e^{-x}\dif x$ diverge si $s\le0$. (Indicación: en $(0,1]$, $e^{-x}\ge e^{-1}$.) Concluya que la integral converge si y solo si $s>0$.
\end{ejercicio}

\begin{ejercicio}{3}{eje:4-limite-comp}
Demuestre el \cref{teo:comparacion-limite}. (Indicación: existe $x_0$ tal que $\frac L2\le\frac{f(x)}{g(x)}\le2L$ para $x\ge x_0$.) Aplíquelo a $\int_1^{\infty}\frac{x+1}{x^{3}+2}\dif x$ y a $\int_1^{\infty}\sin\frac1x\dif x$.
\end{ejercicio}

\begin{ejercicio}{3}{eje:4-dirichlet}
Demuestre que $\int_1^{\infty}\frac{\sin x}{x}\dif x$ converge. (Indicación: integre por partes en $[1,b]$ con $u=\frac1x$ y compare el resultado con el \cref{ej:criterios}(b).) Deduzca que $\int_0^{\infty}\frac{\sin x}{x}\dif x$ converge, y explique por qué en $x=0$ no hay singularidad.
\end{ejercicio}

\begin{ejercicio}{3}{eje:4-potencias}
Sea $\alpha\in\R$ y $f(x)=(\alpha-1)\,x^{-\alpha}$ para $x\ge1$. (a) Demuestre que $f$ es una densidad en $[1,\infty)$, es decir, $f\ge0$ e $\int_1^{\infty}f=1$, si y solo si $\alpha>1$. (b) Para $\alpha>1$, demuestre que la media $\int_1^{\infty}x\,f(x)\dif x$ es finita si y solo si $\alpha>2$, y que entonces vale $\frac{\alpha-1}{\alpha-2}$. (c) ¿Para qué $\alpha$ es finito $\int_1^{\infty}x^{2}f(x)\dif x$?
\end{ejercicio}

\begin{ejercicio}{3}{eje:4-decreciente}
Sea $f\ge0$ decreciente en $[1,\infty)$ con $\int_1^{\infty}f$ convergente. Demuestre que $\lim_{x\to\infty}x\,f(x)=0$. (Indicación: $\int_{x/2}^{x}f\ge\frac x2 f(x)$ para $x\ge2$.) Deduzca otra vez que $\int_1^{\infty}\frac{\dif x}{x}$ diverge.
\end{ejercicio}

\begin{ejercicio}{3}{eje:4-wallis-gauss}
Este ejercicio demuestra el \cref{teo:gauss} sin integrales dobles. Sea $J=\int_0^{\infty}e^{-x^{2}}\dif x$ e $I_n=\int_0^{\pi/2}\sin^{n}x\dif x$ como en el \cref{eje:3-wallis}.
(a) Usando $e^{t}\ge1+t$, demuestre que $(1-x^{2})^{n}\le e^{-nx^{2}}$ en $[0,1]$ y que $e^{-nx^{2}}\le(1+x^{2})^{-n}$ en $[0,\infty)$.
(b) Con los cambios $x=\cos\theta$ y $x=\tan\theta$, deduzca que $\sqrt n\,I_{2n+1}\le J\le\sqrt n\,I_{2n-2}$ para todo $n\ge1$.
(c) De $nI_nI_{n-1}=\frac\pi2$ y del decrecimiento de $(I_n)$, deduzca que $\sqrt{\frac{\pi}{2(n+1)}}\le I_n\le\sqrt{\frac{\pi}{2n}}$, y concluya que $J=\frac{\sqrt\pi}{2}$.
\end{ejercicio}
